\chapter{Introducción}\label{chap:introduccion}

\section{Motivación}

El diagnóstico biológico de la enfermedad de Alzheimer actualmente es clínico y se apoya hoy en biomarcadores sólidos, en especial en los perfiles de amiloide y tau obtenidos en líquido cefalorraquídeo o mediante PET \parencite{jack2018nia}. El reto no es la falta de herramientas, sino su uso como primer filtro, ya que, la punción lumbar es invasiva, la neuroimagen tiene un coste elevado y ambas técnicas presentan una disponibilidad limitada a su vez siendo las más recurridas. Como consecuencia, en la práctica clínica no todos los pacientes con sospecha de deterioro cognitivo pueden estudiarse con la misma intensidad desde el inicio.

En ese contexto, la metabolómica plasmática ofrece una alternativa atractiva: el análisis de sangre es accesible, repetible y potencialmente escalable. Si el perfil metabólico periférico contiene señal suficiente, puede emplearse como primera capa de decisión para priorizar pacientes, descartar casos con un margen razonable de seguridad y reservar pruebas invasivas para quienes realmente las necesitan.

En estudios clínicos con tamaño muestral limitado, el principal riesgo no es la falta de algoritmos, sino la flexibilidad metodológica excesiva: selección de variables inestable, optimismo en validación interna, sobreinterpretación de métricas aisladas y traducción precipitada de probabilidades a decisiones clínicas. Por ello, este problema no se resuelve simplemente entrenando un clasificador sobre una matriz de metabolitos, sino que

desde esta perspectiva, el trabajo no se plantea como una búsqueda del ``mejor'' modelo en términos puramente cuantitativos, sino como un análisis completo del problema, de sus restricciones y de su reformulación más útil en clave clínica.

\section{Planteamiento del problema y contribución del TFM}

El trabajo parte de una formulación clásica de clasificación binaria bajo la premisa de una pregunta inicial directa: si, a partir de un bloque reducido de metabolitos residuales, un modelo supervisado puede discriminar entre AD y NC con una solidez suficiente como para tener interés clínico.

Este punto de partida se apoya en el artículo ``Machine Learning Approach to Select Small Compounds in Plasma as Predictors of Alzheimer’s Disease'' \parencite{bernal2025metabolomics}, donde se muestra que una selección de metabolitos plasmáticos residualizados respecto a un metabolito concomitante de su misma familia biológica puede retener señal predictiva de EA.

Tomando como referencia el dataset residualizado de 20 metabolitos empleado por Bernal et al.~\parencite{bernal2025metabolomics}, para cada paciente se dispone de un vector de entrada $\mathbf{x} \in \mathbb{R}^p$ con información metabolómica residualizada y una etiqueta diagnóstica $y \in \{0,1\}$ que distingue entre control cognitivamente normal (NC) y enfermedad de Alzheimer (AD). En esta formulación inicial, el objetivo es estimar $P(y=1\mid\mathbf{x})$ con capacidad de generalización suficiente para apoyar un triaje preliminar, no para sustituir a las técnicas confirmatorias.

Esa formulación es necesaria, pero no suficiente para tenerla en cuenta de manera directa en un entorno clínico, sin embargo, a medida que avanza el análisis aparecen tres hechos que modifican la naturaleza del problema. Primero, la señal existe, pero con solapamiento relevante entre clases y un margen inevitable de ambigüedad. Segundo, los falsos negativos tienen un coste clínico mayor que los falsos positivos: no detectar a un paciente con AD es más grave que derivar a evaluación adicional a un control sano. Tercero, las probabilidades del modelo no deben leerse solo como salida predictiva, sino como una medida de confianza útil para la toma de decisiones clínicas.

En ese marco, la contribución central del trabajo consiste en recorrer ese cambio de planteamiento de forma explícita y defendible. El estudio comienza con la construcción de un clasificador metabolómico estable y concluye que \textbf{la aportación no es maximizar BA, sino demostrar que la incertidumbre debe modelarse explícitamente} para manejarla y poder tomar decisiones clínicas informadas.

En resumen, el mayor valor metodológico no está en la clasificación binaria aislada, sino en la \textbf{reformulación del problema como un sistema de triaje clínico en dos etapas, con activación condicional del segundo escalón y gestión explícita de la incertidumbre}.

Desde esa perspectiva, los objetivos del trabajo son los siguientes:

\begin{enumerate}
    \item \textbf{Caracterizar rigurosamente el dataset} y determinar qué restricciones impone al modelado: solapamiento entre clases, dependencia entre variables, presencia de extremos y limitaciones de tamaño muestral.
    \item \textbf{Construir un modelo base robusto} mediante selección de variables estable, ingeniería de variables y comparación de distintos modelos de aprendizaje automático, respaldados por criterios estadísticos y clínicos.
    \item \textbf{Evaluar críticamente las limitaciones del clasificador binario}, con especial atención a los falsos negativos y a la ambigüedad biológica de ciertos pacientes.
    \item \textbf{Reformular el problema como triaje}, introduciendo umbrales de confianza y una arquitectura de dos etapas que combine información metabolómica y clínica.
    \item \textbf{Validar el sistema completo de forma robusta}, comparándolo con formulaciones alternativas y verificando qué mejora realmente, qué sacrifica y por qué la solución final no responde a una elección oportunista.
    \item \textbf{Explorar la aplicabilidad potencial de la solución completa}, como herramienta de priorización diagnóstica previa, delimitando con claridad su alcance frente a los métodos confirmatorios de referencia.
\end{enumerate}

\section{Estructura del documento}

La organización del documento sigue la lógica del propio proceso de investigación. Cada capítulo responde a una pregunta concreta y abre el paso al siguiente.

\begin{itemize}
    \item El \textbf{Capítulo~\ref{chap:estado_arte}} sitúa el trabajo en el contexto clínico y metodológico adecuado: biomarcadores de Alzheimer, metabolómica plasmática, aprendizaje automático en muestras pequeñas y sistemas de apoyo a la decisión clínica.

    \item El \textbf{Capítulo~\ref{chap:materiales}} desarrolla el recorrido metodológico completo: datos, análisis exploratorio, selección de variables, construcción del modelo base, análisis de sus límites, reformulación como triaje y diseño formal del sistema.

    \item El \textbf{Capítulo~\ref{chap:resultados}} presenta los hallazgos siguiendo la misma progresión narrativa que estructura el Capítulo 3: qué mostró el EDA, qué rendimiento alcanzó el clasificador base, qué límites reveló el análisis de errores, qué aporta el sistema de dos etapas y cuán robusta es la solución final.

    \item El \textbf{Capítulo~\ref{chap:conclusiones}} sintetiza la aportación del trabajo, discute sus limitaciones y plantea las líneas de continuidad más relevantes.
\end{itemize}
